% ============================================================
% Paper 95 (DE): Uniformisierung in höheren Dimensionen:
%                Yaus Vermutung und Komplexe Geometrie
% Autor: Michael Fuhrmann
% Projekt: Specialist Mathematics Project
% Build: 169
% Datum: 2026-03-12
% ============================================================
\documentclass[12pt,a4paper]{amsart}

\usepackage{amsmath,amssymb,amsthm,mathtools}
\usepackage{enumitem}
\usepackage{booktabs}
\usepackage{array}
\usepackage[hidelinks]{hyperref}
\usepackage[utf8]{inputenc}
\usepackage[T1]{fontenc}
\usepackage[ngerman]{babel}
\usepackage{geometry}
\geometry{margin=2.5cm}

% ---- Theorem-Umgebungen -------------------------------------
\newtheorem{theorem}{Theorem}[section]
\newtheorem{lemma}[theorem]{Lemma}
\newtheorem{corollary}[theorem]{Korollar}
\newtheorem{proposition}[theorem]{Proposition}
\theoremstyle{definition}
\newtheorem{definition}[theorem]{Definition}
\newtheorem{example}[theorem]{Beispiel}
\theoremstyle{remark}
\newtheorem{remark}[theorem]{Bemerkung}
\newtheorem{conjecture}[theorem]{Vermutung}
% Deutsche Aliasse
\newtheorem{satz}[theorem]{Satz}
\newtheorem{korollar}[theorem]{Korollar}
\newtheorem{vermutung}[theorem]{Vermutung}
\newtheorem{bemerkung}[theorem]{Bemerkung}
\newtheorem{definition_de}[theorem]{Definition}

% ---- Operatoren ---------------------------------------------
\DeclareMathOperator{\Ric}{Ric}
\DeclareMathOperator{\Bsc}{Bsc}
\DeclareMathOperator{\vol}{vol}
\DeclareMathOperator{\diam}{diam}
\newcommand{\CC}{\mathbb{C}}
\newcommand{\RR}{\mathbb{R}}
\newcommand{\ZZ}{\mathbb{Z}}
\newcommand{\PP}{\mathbb{P}}
\newcommand{\TT}{\mathbb{T}}
\newcommand{\dd}{\partial}

% ---- Zusammenfassung ----------------------------------------
\newcommand{\OO}{\mathcal{O}}
\begin{document}
\begin{abstract}
Der klassische Uniformisierungssatz klassifiziert alle Riemannschen Flächen
bis auf Biholomorphismus.  In höheren Dimensionen liefern die Sätze von
Mori (kompakter Fall, 1979) und Siu--Yau (kompakte Kähler-Mannigfaltigkeit,
1980) analoge Klassifikationsergebnisse: eine kompakte Kähler-Mannigfaltigkeit
mit positiver bisektionaler Krümmung ist biholomorph zu $\PP^n$.
Yaus Vermutung---dass eine \emph{vollständige} (nicht-kompakte) Kähler-Mannig\-faltigkeit
mit positiver bisektionaler Krümmung biholomorph zu $\CC^n$ ist---bleibt
in voller Allgemeinheit offen.  Wir behandeln Caos Kähler--Ricci-Fluss-Ansatz,
Pseudokonvexität, Stein-Mannigfaltigkeiten und verwandte offene Probleme.
\end{abstract}

\title{Uniformisierung in höheren Dimensionen:\\
\large Yaus Vermutung und Komplexe Geometrie}

\author{Michael Fuhrmann}
\address{Specialist Mathematics Project, Build 169}
\email{specialist-maths@project.local}
\date{2026-03-12}

\maketitle
\tableofcontents

% =============================================================
\section{Einleitung}
% =============================================================

Der \emph{Uniformisierungssatz} für Riemannsche Flächen, bewiesen von
Poincaré und Koebe, ist einer der Grundpfeiler der komplexen Analysis:

\begin{satz}[Poincaré--Koebe Uniformisierung, 1907]
Jede einfach zusammenhängende Riemannsche Fläche ist biholomorph zu genau
einem der Folgenden:
\begin{enumerate}[label=(\roman*)]
  \item der Riemannschen Sphäre $\hat{\CC} \cong S^2$,
  \item der komplexen Ebene $\CC$,
  \item der oberen Halbebene $\mathbb{H} = \{z \in \CC : \Im z > 0\}$.
\end{enumerate}
\end{satz}

Das natürliche höherdimensionale Analogon fragt: Wann ist eine vollständige
Kähler-Mannigfaltigkeit mit positiver Krümmung biholomorph zu $\CC^n$ oder
$\PP^n$?  Yau formulierte dies in den 1970er Jahren als Vermutung,
und der kompakte Fall wurde vollständig gelöst.  Der nicht-kompakte Fall
bleibt eines der zentralen offenen Probleme der komplexen Geometrie.

% =============================================================
\section{Kähler-Mannigfaltigkeiten und bisektionale Krümmung}
% =============================================================

\begin{definition_de}[Kähler-Mannigfaltigkeit]
Eine komplexe Mannigfaltigkeit $(M, J)$ der komplexen Dimension $n$ ist
\emph{kählersch}, wenn sie eine hermitesche Metrik $g$ besitzt, deren
zugehörige $(1,1)$-Form $\omega(X,Y) = g(JX, Y)$ geschlossen ist:
$d\omega = 0$.  Äquivalent: der Levi-Civita-Zusammenhang ist $J$-invariant:
$\nabla J = 0$.
\end{definition_de}

\begin{example}
\begin{enumerate}[label=(\roman*)]
  \item $\CC^n$ mit der euklidischen Metrik (flache Kähler-Mannigfaltigkeit).
  \item $\PP^n$ mit der Fubini--Study-Metrik (positive Schnittkrümmung).
  \item Komplexe Tori $\CC^n/\Lambda$ (flach kählersch).
  \item Glatte projektive Varietäten (Kodaira-Einbettung + Fubini--Study).
\end{enumerate}
\end{example}

\begin{definition_de}[Bisektionale Krümmung \cite{GoldbergKobayashi1967}]
Sei $(M, g, J)$ eine Kähler-Mannigfaltigkeit mit Krümmungstensor $R$.
Für Einheitstangentialvektoren $X, Y \in T_p M$ ist die
\emph{holomorphe bisektionale Krümmung}
\[
  \Bsc(X,Y) = R(X, JX, Y, JY).
\]
Die Mannigfaltigkeit hat \emph{positive bisektionale Krümmung}, wenn
$\Bsc(X,Y) > 0$ für alle nicht-null $X, Y \in T_p M$ und alle $p \in M$ gilt.
\end{definition_de}

\begin{bemerkung}
Positive bisektionale Krümmung impliziert positive Schnittkrümmung,
die wiederum positive Ricci-Krümmung impliziert.
Die Implikationen sind echt: $\CC^n$ hat verschwindende bisektionale Krümmung;
die Fubini--Study-Metrik auf $\PP^n$ hat konstante positive bisektionale
Krümmung.
\end{bemerkung}

% =============================================================
\section{Der klassische Uniformisierungssatz}
% =============================================================

\begin{proof}[Beweisstruktur des Uniformisierungssatzes]
\textbf{Schritt 1 (Harmonische Abbildungen).}  Für eine Riemannsche Fläche
$(\Sigma, g)$ konstruiert man eine harmonische Funktion (Lösung eines
Dirichlet-Problems), um eine konforme Abbildung auf eines der drei Modelle
zu erzeugen.

\textbf{Schritt 2 (Krümmungscharakterisierung).}  Die drei Modelle
entsprechen konstanter Krümmung $+1$ (Sphäre), $0$ (Ebene), $-1$ (Halbebene).
Nach dem Gauß--Bonnet-Satz bestimmt die Eulercharakteristik, welches Modell
zutrifft. $\square$
\end{proof}

% =============================================================
\section{Moris Bend-and-Break-Satz (kompakter Fall)}
% =============================================================

\begin{satz}[Mori, 1979 \cite{Mori1979}]
\label{satz:mori}
Jede glatte kompakte Kähler-Mannigfaltigkeit mit positiver bisektionaler
Krümmung ist biholomorph zu $\PP^n$.
\end{satz}

\begin{satz}[Bend-and-Break \cite{Mori1979}]
Sei $M$ eine glatte projektive Varietät der Dimension $n$ über einem
algebraisch abgeschlossenen Körper.  Enthält $M$ eine glatte Kurve $C$ mit
$(-K_M \cdot C) > 0$, dann enthält $M$ eine rationale Kurve.
\end{satz}

\begin{proof}[Beweisidee für Satz~\ref{satz:mori}]
\textbf{Schritt 1.}  Positive bisektionale Krümmung impliziert $c_1(M) > 0$,
also ist $-K_M$ ampel.

\textbf{Schritt 2.}  Mit Charakteristik-$p$-Methoden (Frobenius) deformiert
Mori Kurven in $M$ (Bend-and-Break) und erzeugt rationale Kurven,
die $M$ überdecken.

\textbf{Schritt 3.}  Nach Kobayashi--Ochiai gilt: Falls $M$ den Index
$i(M) = n+1$ hat, ist $M \cong \PP^n$. $\square$
\end{proof}

% =============================================================
\section{Siu--Yaus Satz: Der differentialgeometrische Beweis}
% =============================================================

\begin{satz}[Siu--Yau, 1980 \cite{SiuYau1980}]
\label{satz:siuyau}
Jede kompakte Kähler-Mannigfaltigkeit mit positiver bisektionaler Krümmung
ist biholomorph zu $\PP^n$.
\end{satz}

\begin{proof}[Beweisidee]
\textbf{Schritt 1 (Frankels Vermutung).}  Die Aussage wurde 1961 von Frankel
vermutet als: "Eine kompakte Kähler-Mannigfaltigkeit mit positiver bisektionaler
Krümmung ist homöomorph zu $\PP^n$."

\textbf{Schritt 2 (Harmonische Abbildungen).}  Mittels harmonischer Abbildungen
und Bochner--Weitzenböck-Formeln zeigen Siu--Yau die Existenz einer
holomorphen Abbildung $f \colon M \to \PP^n$ vom Grad 1.

\textbf{Schritt 3 (Biholomorphismus).}  Die Grad-1-Abbildung $f$ ist ein
Biholomorphismus. $\square$
\end{proof}

\begin{bemerkung}
Die Hypothese \emph{positiver} (nicht nur nicht-negativer) bisektionaler
Krümmung ist wesentlich.  Nicht-negative bisektionale Krümmung führt zu einer
schwächeren Schlussfolgerung (Moks Satz, 1988: biholomorph zu einem kompakten
hermiteschen symmetrischen Raum).
\end{bemerkung}

% =============================================================
\section{Yaus Vermutung: Der nicht-kompakte Fall}
% =============================================================

\begin{vermutung}[Yau, 1970er, \cite{Yau1977} -- OFFEN in voller Allgemeinheit]
\label{verm:yau}
Sei $(M, g)$ eine vollständige Kähler-Mannigfaltigkeit mit positiver
bisektionaler Krümmung.  Dann ist $M$ biholomorph zu $\CC^n$.
\end{vermutung}

\begin{bemerkung}
"Vollständig" bedeutet geodätisch vollständig (jede Geodäte lässt sich
zu einem vollständigen Weg auf $\RR$ fortsetzen).
\end{bemerkung}

\begin{satz}[Teilresultate zu Vermutung~\ref{verm:yau}]
\begin{enumerate}[label=(\roman*)]
  \item \textbf{Komplexe Dimension 2 (Mok, 1984 \cite{Mok1984})}:
        Eine vollständige Kähler-Fläche mit positiver bisektionaler Krümmung
        und endlichem Volumen ist biholomorph zu $\CC^2$.
  \item \textbf{Beschränkte Krümmung (Shi, 1989 \cite{Shi1989})}:
        Mit zusätzlich beschränkter Krümmung gilt $M \cong \CC^n$.
  \item \textbf{Maximales Volumenwachstum (Chen--Tang--Zhu, 2004
        \cite{ChenTangZhu2004})}: Mit $\vol(B(p,r)) \ge c r^{2n}$
        gilt $M \cong \CC^n$.
  \item \textbf{Maximales Volumenwachstum — vollständiger Beweis (Gang Liu, 2016
        \cite{GangLiu2016})}: Gang Liu gab einen vollständigen,
        in sich geschlossenen Beweis, dass jede vollständige Kähler-Mannigfaltigkeit
        mit nicht-negativer bisektionaler Krümmung und \emph{maximalem
        Volumenwachstum} biholomorph zu $\CC^n$ ist.  Dies beseitigt die
        Beschränktheitshypothese von Shi und vereinfacht das Argument
        von Chen--Tang--Zhu.
\end{enumerate}
\end{satz}

\begin{bemerkung}
Lius Satz von 2016 \cite{GangLiu2016} löst Yaus Vermutung im Fall
\emph{maximalen Volumenwachstums} vollständig ohne Krümmungsbeschränktheit.
Der verbleibende offene Fall ist \emph{sub-maximales} Volumenwachstum
(d.\,h.\ $\vol(B(p,r)) = o(r^{2n})$ für $r \to \infty$) kombiniert mit
positiver (nicht nur nicht-negativer) bisektionaler Krümmung.
\end{bemerkung}

% =============================================================
\section{Der Kähler--Ricci-Fluss-Ansatz}
% =============================================================

Hamiltons Ricci-Fluss \cite{Hamilton1982} deformiert eine Riemannsche Metrik
durch $\frac{\partial g}{\partial t} = -2\Ric(g)$.  Für Kähler-Mannigfaltigkeiten
zeigte Cao \cite{Cao1985}, dass der Fluss die Kähler-Eigenschaft erhält.

\begin{satz}[Cao, 1985 \cite{Cao1985}]
Der Kähler--Ricci-Fluss
\[
  \frac{\partial g_{i\bar{j}}}{\partial t} = -R_{i\bar{j}}
\]
hat eine eindeutige Lösung für alle $t \in [0,\infty)$, ausgehend von jeder
kompakten Kähler-Metrik mit positiver bisektionaler Krümmung.  Der Fluss
konvergiert (nach Normierung) zur Fubini--Study-Metrik auf $\PP^n$.
\end{satz}

\begin{satz}[Hamilton, 1988 \cite{Hamilton1988}]
In komplexer Dimension 1 konvergiert der normierte Kähler--Ricci-Fluss zu
einer Metrik konstanter Krümmung.  Dies liefert einen analytischen Beweis
des Uniformisierungssatzes für kompakte Riemannsche Flächen.
\end{satz}

Strategie für Yaus Vermutung via Kähler--Ricci-Fluss:
\begin{enumerate}[label=(\roman*)]
  \item Starte Kähler--Ricci-Fluss auf $(M,g)$ mit positiver bisektionaler Krümmung.
  \item Zeige, dass der Fluss für alle Zeiten existiert.
  \item Zeige, dass der Fluss zur flachen Metrik auf $\CC^n$ konvergiert.
  \item Schließe $M \cong \CC^n$.
\end{enumerate}
Schritte (i) und (ii) sind unter Krümmungsschranken bekannt;
Schritte (iii) und (iv) erfordern weitere Arbeit.

% =============================================================
\section{Stein-Mannigfaltigkeiten und Pseudokonvexität}
% =============================================================

\begin{definition_de}[Stein-Mannigfaltigkeit]
Eine komplexe Mannigfaltigkeit $M$ ist \emph{Stein}, wenn sie erfüllt:
\begin{enumerate}[label=(\roman*)]
  \item \textbf{Holomorphe Konvexität}: Für jedes kompakte $K \subseteq M$
        ist die holomorphe Hülle $\hat K_\OO$ kompakt.
  \item \textbf{Holomorphe Trennbarkeit}: Für $p \neq q$ gibt es eine
        holomorphe Funktion $f$ mit $f(p) \neq f(q)$.
  \item \textbf{Lokale holomorphe Koordinaten}: An jedem Punkt spannen
        die Differentiale holomorpher Funktionen den Kotangentialraum auf.
\end{enumerate}
\end{definition_de}

\begin{satz}[Grauert, 1958 \cite{Grauert1958}]
Eine komplexe Mannigfaltigkeit $M$ ist genau dann Stein, wenn sie eine
streng plurisubharmonische Erschöpfungsfunktion
$\varphi \colon M \to \RR$ besitzt ($i\dd\bar\dd\varphi > 0$ im Sinne
von $(1,1)$-Formen, mit kompakten Unterniveaumengen).
\end{satz}

\begin{satz}[Greene--Wu, 1978 \cite{GreeneWu1978}]
Eine vollständige Kähler-Mannigfaltigkeit mit positiver bisektionaler
Krümmung ist Stein.
\end{satz}

\begin{korollar}
Yaus Vermutung würde bedeuten: Die einzige Steinschen Kähler-Mannigfaltigkeit
mit positiver bisektionaler Krümmung ist $\CC^n$.
\end{korollar}

% =============================================================
\section{Die Hyperkonvexitäts-Vermutung}
% =============================================================

\begin{definition_de}[Hyperkonvexes Gebiet]
Ein Gebiet $\Omega \subseteq \CC^n$ heißt \emph{hyperkonvex}, wenn es eine
beschränkte plurisubharmonische Erschöpfungsfunktion
$u \colon \Omega \to [-\infty, 0)$ besitzt.
\end{definition_de}

\begin{vermutung}[Hyperkonvexitäts-Vermutung -- OFFEN]
\label{verm:hyperkonvex}
Jedes pseudokonvexe Gebiet in $\CC^n$ mit $C^1$-Rand ist hyperkonvex.
\end{vermutung}

\begin{satz}[Diederich--Fornæss, 1977 \cite{DiederichFornaess1977}]
Jedes beschränkte pseudokonvexe Gebiet mit $C^2$-Rand in $\CC^n$ ist hyperkonvex.
\end{satz}

Die Vermutung für $C^1$-Ränder bleibt offen.
Gegenbeispiele existieren für lediglich Lipschitz-reguläre Ränder
(Wurm-Gebiete nach Diederich--Fornæss).

% =============================================================
\section{Moks Satz über nicht-negative bisektionale Krümmung}
% =============================================================

\begin{satz}[Mok, 1988 \cite{Mok1988}]
\label{satz:mok}
Sei $(M, g)$ eine kompakte irreduzible Kähler-Mannigfaltigkeit mit
nicht-negativer bisektionaler Krümmung.  Dann ist $M$ biholomorph zu
einem der folgenden \emph{kompakten irreduziblen hermiteschen symmetrischen Räume}:
\begin{enumerate}[label=(\roman*)]
  \item $\PP^n$ (komplexer projektiver Raum),
  \item $\mathrm{Gr}(k,n)$ (komplexe Grassmannsche Varietät),
  \item $\mathrm{SO}(n)/U(n/2)$ (Typ DIII),
  \item $\mathrm{Sp}(n)/U(n)$ (Typ CI),
  \item $E_6/\mathrm{Spin}(10)\cdot U(1)$,
  \item $E_7/E_6 \cdot U(1)$.
\end{enumerate}
\end{satz}

\begin{bemerkung}
Moks Satz ist die vollständige Klassifikation unter der schwächeren Hypothese
nicht-negativer (statt strikt positiver) bisektionaler Krümmung.
\end{bemerkung}

% =============================================================
\section{Vergleich: Riemannsche Flächen vs. höhere Dimensionen}
% =============================================================

\begin{center}
\begin{tabular}{p{2.8cm}p{3.5cm}p{3.5cm}p{3.5cm}}
\toprule
 & Riemannsche Fläche & Kompakt kählersch & Vollständig kählersch \\
\midrule
Pos. Krümmung & $S^2 \cong \PP^1$ & $M \cong \PP^n$ (Mori--Siu--Yau, bewiesen) & $M \cong \CC^n$? (Yau, offen) \\[4pt]
Null-Krümmung & $\CC$ oder $\TT^2$ & Kähler-flach (CY) & Flach ($\CC^n/\Gamma$) \\[4pt]
Neg. Krümmung & $\mathbb{H}$ & Allgemeiner Typ & Cartan--Hadamard \\
\bottomrule
\end{tabular}
\end{center}

% =============================================================
\section{Offene Probleme}
% =============================================================

\begin{vermutung}[Yaus Vermutung über positive bisektionale Krümmung,
= Vermutung~\ref{verm:yau} -- OFFEN in voller Allgemeinheit]
Jede vollständige Kähler-Mannigfaltigkeit mit positiver bisektionaler
Krümmung ist biholomorph zu $\CC^n$.
\end{vermutung}

\begin{vermutung}[Yaus Uniformisierungsvermutung für positive Ricci-Krümmung -- OFFEN]
Jede vollständige Kähler-Mannigfaltigkeit mit positiver Ricci-Krümmung,
beschränkter Krümmung und maximalem Volumenwachstum ist biholomorph zu $\CC^n$.
\end{vermutung}

\begin{vermutung}[Frankelsche Vermutung (nicht-kompakte Version) -- OFFEN]
Jede vollständige einfach zusammenhängende Kähler-Mannigfaltigkeit mit
positiver Schnittkrümmung ist biholomorph zu $\CC^n$.
\end{vermutung}

\begin{vermutung}[Hyperkonvexität für $C^1$-Ränder,
= Vermutung~\ref{verm:hyperkonvex} -- OFFEN]
Jedes pseudokonvexe Gebiet in $\CC^n$ mit $C^1$-Rand ist hyperkonvex.
\end{vermutung}

% =============================================================
\section{Fazit}
% =============================================================

Der Uniformisierungssatz für Riemannsche Flächen hat Jahrzehnte des Bemühens
angeregt, komplexe Mannigfaltigkeiten in höheren Dimensionen zu verstehen.
Der kompakte Fall von Yaus Vermutung wurde von Mori und Siu--Yau gelöst:
Kompakte Kähler-Mannigfaltigkeiten mit positiver bisektionaler Krümmung
sind $\PP^n$.

Der nicht-kompakte Fall---Yaus Vermutung im eigentlichen Sinne---bleibt in
voller Allgemeinheit offen.  Teilresultate unter Krümmungsschranken (Shi, 1989)
und maximalem Volumenwachstum (Chen--Tang--Zhu, 2004) stützen die Vermutung.
Der Kähler--Ricci-Fluss-Ansatz (Cao, 1985) liefert ein potenziell mächtiges
Werkzeug, aber das Langzeit-Verhalten des Flusses im nicht-kompakten Fall
ohne beschränkte Krümmung ist noch nicht vollständig verstanden.

\begin{thebibliography}{99}

\bibitem{Yau1977}
S.-T. Yau,
\textit{Calabi's conjecture and some new results in algebraic geometry},
Proc. Nat. Acad. Sci. USA \textbf{74} (1977), 1798--1799.

\bibitem{Mori1979}
S.~Mori,
\textit{Projective manifolds with ample tangent bundles},
Ann. of Math. (2) \textbf{110} (1979), 593--606.

\bibitem{SiuYau1980}
Y.-T. Siu und S.-T. Yau,
\textit{Compact Kähler manifolds of positive bisectional curvature},
Invent. Math. \textbf{59} (1980), 189--204.

\bibitem{Mok1984}
N.~Mok,
\textit{An embedding theorem of complete Kähler manifolds of positive bisectional curvature onto affine algebraic varieties},
Bull. Soc. Math. France \textbf{112} (1984), 197--258.

\bibitem{Mok1988}
N.~Mok,
\textit{The uniformization theorem for compact Kähler manifolds of nonneg\-ative holomorphic bisectional curvature},
J. Differential Geom. \textbf{27} (1988), 179--214.

\bibitem{Cao1985}
H.-D. Cao,
\textit{Deformation of Kähler metrics to Kähler--Einstein metrics on compact
Kähler manifolds},
Invent. Math. \textbf{81} (1985), 359--372.

\bibitem{Hamilton1982}
R.~S. Hamilton,
\textit{Three-manifolds with positive Ricci curvature},
J. Differential Geom. \textbf{17} (1982), 255--306.

\bibitem{Hamilton1988}
R.~S. Hamilton,
\textit{The Ricci flow on surfaces},
Contemp. Math. \textbf{71} (1988), 237--262.

\bibitem{Shi1989}
W.-X. Shi,
\textit{Ricci deformation of the metric on complete noncompact Riemannian
manifolds},
J. Differential Geom. \textbf{30} (1989), 223--301.

\bibitem{ChenTangZhu2004}
B.-L. Chen, S.-H. Tang und X.-P. Zhu,
\textit{A uniformization theorem for complete non-compact Kähler surfaces
with positive bisectional curvature},
J. Differential Geom. \textbf{67} (2004), 519--570.

\bibitem{GoldbergKobayashi1967}
S.~I. Goldberg und S.~Kobayashi,
\textit{Holomorphic bisectional curvature},
J. Differential Geom. \textbf{1} (1967), 225--233.

\bibitem{Grauert1958}
H.~Grauert,
\textit{On Levi's problem and the imbedding of real-analytic manifolds},
Ann. of Math. (2) \textbf{68} (1958), 460--472.

\bibitem{GreeneWu1978}
R.~E. Greene und H.~Wu,
\textit{$C^\infty$ approximations of convex, subharmonic, and plurisubharmonic
functions},
Ann. Sci. Éc. Norm. Supér. (4) \textbf{12} (1979), 47--84.

\bibitem{GangLiu2016}
G.~Liu,
\textit{Three-circle theorem and dimension estimate for holomorphic functions on K\"ahler manifolds},
Duke Math.\ J.\ \textbf{165} (2016), Nr.~15, 2899--2919.

\bibitem{DiederichFornaess1977}
K.~Diederich und J.~E. Fornæss,
\textit{Pseudoconvex domains: bounded strictly plurisubharmonic exhaustion
functions},
Invent. Math. \textbf{39} (1977), 129--141.

\end{thebibliography}

\end{document}
